\begin{topic}[id=sdv_zonal_ee_arch,
             difficulty={Anspruchsvoll},
             type={Architektur- und Plattformanalyse},
             prerequisites={Grundlagen zu eingebetteten Systemen, Rechnerarchitektur, Fahrzeugnetzwerken und IT-Sicherheit},
             practice={optional},
             owner={Dozententeam},
             updated={2026-04-13},
             tags={ Software-defined Vehicle, zonale E/E-Architektur, Fahrzeug-Compute, Virtualisierung, OTA-Updates, AUTOSAR Adaptive, Safety, Security }]{Software-defined Vehicle und zonale E/E-Architekturen}

\TopicDescription{Software-defined Vehicle (SDV) bezeichnet eine Fahrzeugarchitektur, in der Funktionen, Updates und Varianten zunehmend über Software statt über eine große Zahl fest verdrahteter Steuergeräte differenziert werden. Für ein seminarfähiges Thema sollte der Fokus daher nicht auf Marketingbegriffen liegen, sondern auf der technischen Ablösung klassischer, funktionszentrierter ECU-Landschaften durch zonale E/E-Architekturen mit zentralem Fahrzeug-Compute, Ethernet-Backbone und klarer Trennung zwischen hardware-nahen und adaptiven Softwareebenen. Im Mittelpunkt stehen die Rollen von Zonencontrollern, Hochleistungsrechnern, Virtualisierung und service-orientierter Kommunikation sowie die Frage, wie Classic- und Adaptive-AUTOSAR-Anteile gemeinsam betrieben, isoliert und aktualisiert werden. Ebenso relevant ist der Engineering-Wandel: Deployment auf Rechnerplattformen statt nur auf einzelne ECUs, softwaregetriebene Konfiguration, OTA-Updatefähigkeit und Lebenszyklusmanagement über Fahrzeugplattformen hinweg. Damit verbinden sich neue Anforderungen an Safety und Security, etwa sichere Partitionierung, kontrollierte Kommunikation zwischen Zonen, Firewall-Konzepte, Update-Absicherung und Nachweisbarkeit in heterogenen Rechnerplattformen. Auch die Migration ist technisch anspruchsvoll, weil bestehende Funktionen, Bussysteme und Echtzeitanforderungen nicht abrupt verschwinden, sondern in Hybridarchitekturen weitergeführt werden müssen. Das Thema eignet sich deshalb für eine Architektur- und Trade-off-Analyse zwischen Zentralisierung und Restdezentralität, Offenheit und Herstellerbindung, Rechenkonsolidierung und Echtzeitnähe sowie Wiederverwendbarkeit, Wartbarkeit und Integrationsaufwand im künftigen Fahrzeug.}

\TopicLeitfragen{%
\item Welche technischen Ziele verfolgen zonale E/E-Architekturen gegenüber klassischen funktionszentrierten ECU-Landschaften?
\item Wie werden zentrale Fahrzeugrechner, Zonencontroller, Virtualisierung und Classic-/Adaptive-AUTOSAR in einer SDV-Plattform zusammengesetzt?
\item Welche Anforderungen entstehen für Updatefähigkeit, sichere Partitionierung und Lebenszyklusmanagement über mehrere Rechnerdomänen hinweg?
\item Wo liegen Zielkonflikte zwischen Safety, Security, Echtzeit, Rechenkonsolidierung und Engineering-Komplexität?
}

\TopicScope{%
\item Keine Marktübersicht konkreter OEM-, Tier-1- oder Halbleiterstrategien.
\item Keine detaillierte Analyse einzelner Bussysteme wie CAN, LIN oder Automotive Ethernet im Protokolldetail.
\item Kein Fokus auf automatisiertes Fahren, Infotainment oder Cloud-Dienste als eigenständige Themen.
\item Keine vollständige Normenübersicht zu Homologation, Cybersecurity oder funktionaler Sicherheit.
}

\TopicPractice{Anhand einer vereinfachten Fahrzeugarchitektur mit zwei Zonencontrollern, zentralem Rechner und verbliebenen Bestands-ECUs kann ein Migrationspfad vom klassischen Steuergeräteverbund zur zonalen Plattform modelliert werden. Dazu reichen Funktionszuordnung, Kommunikationspfade, Virtualisierungsgrenzen und ein OTA-Update-Szenario ohne Hardware-Aufbau.}

\TopicKeywords{Software-defined Vehicle, zonale E/E-Architektur, Fahrzeug-Compute, Virtualisierung, OTA-Updates, AUTOSAR Adaptive, Safety, Security}

\nocite{veh_sdvzonal_autosar2025,veh_sdvzonal_aparch2023,veh_sdvzonal_apmeth2024,veh_sdvzonal_fw2024,veh_sdvzonal_unr1562021,veh_sdvzonal_sdvall2024}
\end{topic}
